\emph{How the apertures are defined now.}  The metric, the activity rule, the projection and the
VOF sweeps (which freeze a copy at the start of the step) all use the interface-owned phase
apertures of \code{\_current\_phase\_face\_apertures!}, i.e.\ \code{\_phase\_face\_apertures}: the
liquid aperture is the mean of the two cells' traces, $a_f=\tfrac12(t_{c_-}+t_{c_+})$, with the
trace $t_c$ of \code{\_face\_aperture} the part of the face below the PLIC line of $c$ ($0$ for
$\alpha_c\le0$, the full open face for $\alpha_c\ge1$, a solid clipped first), and the gas
aperture $1-a_f$.  The exceptions that clip to zero or one are:
\begin{enumerate}[leftmargin=*,itemsep=1pt,label=(\roman*)]
\item a phase absent in both cells gets aperture exactly $0$ (the other phase the open face);
  this coincides with the plain mean, since both traces are then $0$ or full;
\item the round-off snap \code{\_snap\_phase\_apertures}: an aperture below $64\epsilon$ is set to
  $0$ (and the other to the open face), against $10^{-17}$ openings left by two independently
  clipped polygons;
\item indirectly, the VOF floor: cell fractions within \code{fraction\_floor} ($\sqrt\epsilon$) of
  $0$ or $1$ are snapped after each sweep, so their traces become $0$ or full;
\item in the metric only, a grid-line face between a pure liquid and a pure gas cell is
  identified (one value, each cell sees its own side fully wet), whatever its mean aperture
  ($\tfrac12$);
\item the boundary faces (\code{\_phase\_boundary\_face\_apertures}) use the single trace with
  the same snap.
\end{enumerate}
No pure-neighbour rule is left in this branch (it was removed with MMd, \code{report\_mmd.md}).
The present aperture \emph{is} therefore the plain $\tfrac12$--$\tfrac12$ mean up to (ii)
(round-off) and the grid-line identification (iv).  The prototype switch
\code{\_APERTURE\_PLAIN[]} (environment \code{MM\_APERTURE=plain}) removes (i) and (ii); (iii)
belongs to the VOF and (iv) to the metric and are kept.

\emph{What changes in principle.}  A value is born when either neighbour's trace reaches the face
and dies when neither does --- the same events as now, since (i) is not a clip.  Without (ii) a
phase can be active on a $10^{-17}$ opening: such values have an inertia of order $b\le10^{-16}$
and, under aperture activity, become degrees of freedom; with the small ties their inertia is
that of their staggered mass instead.  The traces move continuously while one PLIC line
translates, but jump when the reconstruction changes its normal between steps (a new Youngs or
LHF normal pivots the line), so births and deaths are not smoother in time with either
definition: the time continuity is set by the reconstruction, not by the averaging.  The Gram
block of a cell uses the shared averaged apertures of its four sides, which need not close a
polygon with any one cell's PLIC (\code{metric\_options}); the VOF volume fluxes are geometric
(donating regions of the PLIC of the upwind cell), not apertures, so neither definition makes the
fluxes and the block's wetted sides consistent.  The identities of \cref{sec:ap:analysis} do not
depend on how the apertures are computed.
